% ---------------------------------------------------------------------------
\section{Evaluation plan}
\label{sec:evaluation}

\todo{This section describes a study that has \emph{not} been run. No result
in this draft comes from participants. Replace with the study once it is
pre-registered and run; until then the core claim of \cref{sec:model}---that
bearing-faithful placement helps---is a design conjecture and must be worded as
one everywhere.}

The model's descriptive and generative power is argued in
\cref{sec:descriptive,sec:generative}; its central design claim is empirical and
untested. Two of its axes make testable predictions, and they interact, so we
plan one controlled, within-subjects study crossing them.

\paragraph{Factors.} \emph{Angle}: categorical (block placement, angle
carries category order) versus bearing (necklace placement, angle carries
bearing). \emph{Anchor}: closed ring versus unrolled axis with leaders. A third,
between-trials factor manipulates \emph{displacement}---the gap between a
mark and its true bearing, at four levels from 0 to 45\,degrees
\todo{levels to be fixed by pilot}---to locate the point at which a displaced
mark misleads (the open policy question of \cref{sec:association}).

\paragraph{Tasks.} Five tasks, each with a ground truth computable from the
stimulus:
\begin{enumerate}[leftmargin=*, itemsep=1pt]
  \item \textbf{Direction}: ``In which direction from the centre is category
    $X$ most concentrated?'' Response: a bearing; error in degrees.
  \item \textbf{Isotropy}: ``Is $X$ spread evenly around the centre?''
    Ground truth from the resultant length $R$.
  \item \textbf{Comparison}: ``Which of two highlighted directions has more
    $X$?'' Accuracy.
  \item \textbf{Composition}: ``Which category is most common inside the
    lens?'' Accuracy. The categorical ring should do at least as well here;
    the task is included so that a win on direction is not bought silently.
  \item \textbf{Association}: ``Click the part of the map that the highlighted
    mark summarises.'' Distance error.
\end{enumerate}

\paragraph{Hypotheses.}
H1: bearing-faithful placement reduces direction error (task~1) relative to
categorical placement.
H2: categorical placement is no worse than bearing placement for composition
(task~4).
H3: the unrolled axis improves comparison accuracy (task~3) relative to the
ring.
H4: the ring beats the unrolled axis on association (task~5) without leaders,
and leaders close the gap.
H5: direction error grows with displacement, and leaders flatten that growth.

\paragraph{Stimuli, participants, analysis.} Synthetic point patterns with
controlled anisotropy (so ground truth is exact) plus a held-out set from real
OpenStreetMap extracts, rendered by the library. Participants recruited online
\todo{platform, $N$ from a power analysis on pilot data, compensation, ethics
approval}. Mixed-effects models with participant and stimulus as random
effects \todo{specify before data collection; pre-register (e.g.\ OSF)}.

\paragraph{Known risk.} A study of off-screen indicators reports that
participants located off-screen objects better, and overwhelmingly preferred,
orthographic rather than radial projection onto the viewport border
\cite{jackle2017topology}. That is evidence \emph{against} direction-preserving
placement on a boundary, in a closely related setting; our hypotheses H1 and
H5 should be read as the direct test of whether it transfers to a lens.
\doubt{Exact result, $N$ and task of Jäckle et al.\ to be confirmed from the
full text before this paragraph is kept.}
Relatedly, a study of orbital boundary labels around a circle reports straight
leaders faster than orbital-radial ones at similar accuracy
\cite{wallinger2026clarity}; our leaders are straight.
\doubt{Pending the citation audit.}
